\section{Seat Allocation Projections}

\subsection{Huntington-Hill on Projected Populations}

We apply the bisect-apportion Huntington-Hill implementation (\texttt{bisect apportion})
to the 2030 medium-series state populations. The 2030 total US population (medium series)
is projected at approximately 340 million. With 435 seats, the ideal district size is
approximately 781,600.

\begin{table}[h]
\centering\small
\caption{Projected 2030 congressional seat changes (medium series, Huntington-Hill).
Confidence column shows range across low/medium/high series.}
\label{tab:2030seats}
\begin{tabular}{lcccc}
\toprule
State & 2020 Seats & 2030 Projected & Change & CI (low--high) \\
\midrule
Texas & 38 & 41 & $+3$ & $+2$ to $+3$ \\
Florida & 28 & 30 & $+2$ & $+2$ to $+3$ \\
Arizona & 9 & 10 & $+1$ & $+1$ to $+1$ \\
Georgia & 14 & 15 & $+1$ & $+1$ to $+1$ \\
Colorado & 8 & 9 & $+1$ & $+0$ to $+1$ \\
North Carolina & 14 & 15 & $+1$ & $+1$ to $+1$ \\
Washington & 10 & 11 & $+1$ & $+1$ to $+1$ \\
Oregon & 6 & 7 & $+1$ & $+0$ to $+1$ \\
\midrule
New York & 26 & 24 & $-2$ & $-3$ to $-2$ \\
California & 52 & 51 & $-1$ & $-1$ to $+0$ \\
Illinois & 17 & 16 & $-1$ & $-1$ to $-1$ \\
Ohio & 15 & 14 & $-1$ & $-1$ to $-1$ \\
Michigan & 13 & 12 & $-1$ & $-1$ to $-1$ \\
Pennsylvania & 17 & 16 & $-1$ & $-1$ to $-1$ \\
West Virginia & 2 & 1 & $-1$ & $-1$ to $-1$ \\
Minnesota & 8 & 7 & $-1$ & $-1$ to $-0$ \\
Rhode Island & 2 & 1 & $-1$ & $-1$ to $-1$ \\
Montana & 2 & 1 & $-1$ & $-1$ to $-0$ \\
\bottomrule
\end{tabular}
\end{table}

\subsection{Algorithmic Implications of Seat Changes}

Each seat change requires designing a new bisection tree.
The ApportionRegions bisection tree depends on the prime factorization of $k$:

\begin{itemize}
  \item \textbf{TX $k = 41$}: 41 is prime. The bisect-apportion binary fallback
    splits 41 into $\lfloor 41/2 \rfloor = 20$ and $\lceil 41/2 \rceil = 21$,
    avoiding a direct 41-way METIS partition. Both 20 and 21 are composite
    ($20 = 4\times5$, $21 = 3\times7$) and handled efficiently.
  \item \textbf{FL $k = 30$}: $30 = 2\times3\times5$. ApportionRegions splits
    30 into two 15-way partitions or three 10-way partitions depending on the
    factorization ordering (smallest-prime-first vs.\ largest-prime-first).
  \item \textbf{GA $k = 15$}: $15 = 3\times5$. Three groups of 5 each.
  \item \textbf{NC $k = 15$}: Same as GA. A significant change from $k=14$
    (which uses $14 = 2\times7$).
  \item \textbf{NY $k = 24$}: $24 = 2^3\times3$. Eight groups of 3 each.
  \item \textbf{WV $k = 1$}: Single district; no redistricting needed.
\end{itemize}

\subsection{Confidence in the Projection}

The medium series is used as the primary projection. The low--high range
represents the 80\% confidence interval around the medium series.
Six states have uncertain seat allocations (marked $\pm1$ in the CI column).
For planning purposes, the bisect team should pre-build bisection trees for
both the baseline and $\pm1$ scenarios for these six states.
